\documentclass[10pt]{article}
\usepackage[margin=0.82in]{geometry}
\usepackage{amsmath,amssymb,amsthm,mathtools}
\usepackage{graphicx}
\usepackage{microtype}
\usepackage[hidelinks]{hyperref}
\usepackage{xcolor}

\newtheorem{theorem}{Theorem}
\newtheorem{lemma}{Lemma}
\newtheorem{proposition}{Proposition}
\newcommand{\C}{\mathbb C}
\newcommand{\N}{\mathbb N}
\newcommand{\Kreiss}{\mathcal K}

\title{Polynomial lower bounds for every fixed Kreiss constant\\
\large Candidate full solution to Question 6.1 of arXiv:2512.10025}
\author{Agent lane 17 (model GPT5.6)}
\date{11 August 2026}

\begin{document}
\maketitle

\begin{center}
\fbox{\parbox{0.92\linewidth}{\textbf{Status.}
Candidate full solution, likely valid; novelty unconfirmed.  The argument
answers Question 6.1 affirmatively.  It does not answer the stronger uniform
exponent Question 6.2.}}
\end{center}

\section{Source question}

The source is N. Chalmoukis, G. Tsikalas, and D. Yakubovich,
\emph{Operators with small Kreiss constants},\\
arXiv:2512.10025v2
\cite{CTY}.  For an $N\times N$ matrix $A$, write
\[
 \Kreiss(A)=\sup_{|z|>1} (|z|-1)\|(zI-A)^{-1}\|.
\]
The source defines $P(N,K)$ as the least universal bound on
$\sup_{n\ge1}\|A^n\|$ over all $N\times N$ matrices with
$\Kreiss(A)\le K$.

Question 6.1 on PDF page 18 asks whether there is a strictly increasing
$\alpha:(1,\infty)\to(0,1)$, with endpoint limits $0$ and $1$, such that
\[
 P(N,K)\gtrsim_K N^{\alpha(K)} \qquad(N\ge1,\ K>1).
\]

\begin{figure}[h]
\centering
\includegraphics[width=0.98\linewidth]{figures/open_problem_crop.png}
\caption{The source context and complete statement of Question 6.1, PDF page 18.}
\end{figure}

\section{Proof intuition}

The source obtains arbitrarily high powers of $\log N$ by increasing an
internal matrix dimension, but its perturbation constants deteriorate too
quickly with that dimension.  The key idea here is to tensorize a single
$2\times2$ nilpotent perturbation instead.  Tensor products amplify a
carefully selected power exponentially in the tensor degree $m$, while a
pointwise comparison collapses the entire perturbation estimate to the scalar
weight $((j+n)/j)^a-1$.  Its averaged square is $O(a^2)$ independently of
$m$.  Abel summation then gives a Kreiss constant $1+O(a)$, whereas choosing
$m\asymp a\log L$ in an $L$-coordinate compression turns the exponential
tensor amplification into a positive power of the total dimension.

\section{Quantitative theorem}

All logarithms are natural.  For $0<a<\tfrac12$, set
\[
 \gamma(a)=a\sqrt{\frac{6}{(1-2a)(1-a)}},
 \qquad
 \beta(a)=\frac{a\log2}{4(1+a\log2)}.
\]

\begin{theorem}\label{thm:quant}
For every $0<a<\tfrac12$ there is $c_a>0$ such that
\[
 P(N,1+\gamma(a))\ge c_a N^{\beta(a)}
 \qquad(N\ge1).
\]
Consequently, Question 6.1 has an affirmative answer.
\end{theorem}

The proof is self-contained except for the previously known large-$K$ lower
bound of Nikolski \cite{Nikolski}, already recorded as equation (1.5) in the
source paper \cite{CTY}.

\section{A dimension-free perturbation estimate}

Let
\[
 J=\begin{pmatrix}0&1\\0&0\end{pmatrix},\qquad J^2=0.
\]
Fix $m\ge1$ and put $H_m=(\C^2)^{\otimes m}$.  On
$\ell^2(\N;H_m)$ define the weighted backward shift $T=T_{a,m}$ by
\[
 (Tx)_j=B_{m,j+1}x_{j+1},\qquad
 B_{m,k}=\left(I_2+\frac{a}{m}\log\frac{k}{k-1}\,J\right)^{\otimes m}
 \quad(k\ge2).
\]
Let $S$ be the unweighted backward shift on the same space.

Because tensor products multiply componentwise and $J^2=0$, the block of
$T^n$ carrying coordinate $j+n$ to coordinate $j$ is exactly
\begin{equation}\label{eq:block}
 C_{j,n}=\left(I_2+\frac{a}{m}\log\frac{j+n}{j}\,J\right)^{\otimes m}.
\end{equation}
For $t\ge0$, telescoping the difference of tensor powers gives
\[
 \|(I_2+tJ)^{\otimes m}-I_{H_m}\|
 \le (1+t)^m-1\le e^{mt}-1.
\]
Hence \eqref{eq:block} implies
\begin{equation}\label{eq:blockdiff}
 \|C_{j,n}-I_{H_m}\|
 \le \left(\frac{j+n}{j}\right)^a-1.
\end{equation}

\begin{lemma}\label{lem:average}
With $D_n=T^n-S^n$, one has, for every $x\in\ell^2(\N;H_m)$ and $M\ge1$,
\[
 \frac1M\sum_{n=1}^M\|D_nx\|\le\gamma(a)\|x\|.
\]
The bound is independent of $m$.
\end{lemma}

\begin{proof}
Put $f_a(u)=(u^{-a}-1)^2$.  A direct integration gives
\begin{equation}\label{eq:integral}
 I_a:=\int_0^1 f_a(u)\,du
 =\frac{2a^2}{(1-2a)(1-a)}.
\end{equation}
By \eqref{eq:blockdiff}, changing variables from $(j,n)$ to $k=j+n$ yields
\[
 \frac1M\sum_{n=1}^M\|D_nx\|^2
 \le\sum_{k\ge2}q_{M,k}\|x_k\|^2,
\]
where
\[
 q_{M,k}=\frac1M\sum_{n=1}^{\min(M,k-1)}
 \left[\left(\frac{k}{k-n}\right)^a-1\right]^2.
\]
We claim $q_{M,k}\le3I_a$.  If $k\le2M+1$, monotonicity of $f_a$ and the
upper Riemann-sum estimate give
\[
 q_{M,k}\le\frac1M\sum_{j=1}^{k-1}f_a(j/k)
 \le\frac{k}{M}I_a\le3I_a.
\]
If $k>2M+1$, then $1-n/k>1/2$ for $1\le n\le M$.  The mean value theorem
gives
\[
 (1-n/k)^{-a}-1\le a2^{a+1}\frac nk,
\]
and therefore
\[
 q_{M,k}\le a^2 2^{2a+2}\left(\frac Mk\right)^2
 \le a^2 2^{2a}<2a^2\le3I_a.
\]
Thus the quadratic average is at most $3I_a\|x\|^2$.  Cauchy--Schwarz in
$n$ and \eqref{eq:integral} finish the proof because
$\sqrt{3I_a}=\gamma(a)$.
\end{proof}

\begin{proposition}\label{prop:kreiss}
$\Kreiss(T_{a,m})\le1+\gamma(a)$ for every $m\ge1$.
\end{proposition}

\begin{proof}
The block formula gives $\|T^n\|\le(n+1)^a$, so the Neumann resolvent series
converges in norm for $|z|>1$.  Let $r=|z|$, $q=r^{-1}$, and
$d_n=\|D_nx\|$.  Lemma \ref{lem:average} says that
$A_N:=\sum_{n=1}^Nd_n\le\gamma(a)N\|x\|$.  Abel summation gives
\begin{align*}
 (r-1)\sum_{n\ge1}r^{-n-1}d_n
 &=(1-q)\sum_{n\ge1}q^nd_n\\
 &=(1-q)^2\sum_{N\ge1}q^N A_N
 \le\gamma(a)q\|x\|\le\gamma(a)\|x\|.
\end{align*}
Since $S$ is a contraction,
\[
 (r-1)\|(zI-S)^{-1}\|
 \le(r-1)\sum_{n\ge0}r^{-n-1}=1.
\]
Writing the resolvent of $T$ as the resolvent series for $S$ plus the series
of the $D_n$ proves the assertion.
\end{proof}

\section{Finite compressions and polynomial growth}

Let $A_{m,L}$ be the restriction of $T_{a,m}$ to its invariant first $L$
coordinates.  It acts on a space of dimension $2^mL$, is nilpotent, and by
Proposition \ref{prop:kreiss} satisfies
\[
 \Kreiss(A_{m,L})\le1+\gamma(a).
\]
Apply $A_{m,L}^{L-1}$ to the unit vector supported at coordinate $L$ with
internal component $e_2^{\otimes m}$.  Formula \eqref{eq:block} gives
\begin{equation}\label{eq:power}
 \|A_{m,L}^{L-1}\|
 \ge\left(1+\left(\frac{a\log L}{m}\right)^2\right)^{m/2}.
\end{equation}
If $a\log L\ge2$ and $m=\lfloor a\log L\rfloor$, then
$m\ge a\log L/2$ and $a\log L/m\ge1$.  Therefore
\begin{equation}\label{eq:Lgrowth}
 \|A_{m,L}^{L-1}\|\ge2^{m/2}\ge L^{a\log2/4},
 \qquad
 \dim A_{m,L}=2^mL\le L^{1+a\log2}.
\end{equation}

Given a sufficiently large target dimension $N$, take
\[
 L=\left\lfloor N^{1/(1+a\log2)}\right\rfloor,
 \qquad m=\lfloor a\log L\rfloor.
\]
Then $2^mL\le N$.  Pad $A_{m,L}$ with a zero block to obtain an $N\times N$
matrix without increasing its Kreiss constant.  The floor losses are absorbed
in a constant depending only on $a$, and \eqref{eq:Lgrowth} yields
\[
 P(N,1+\gamma(a))\gtrsim_a
 N^{\frac{a\log2}{4(1+a\log2)}}.
\]
The finitely many smaller $N$ are absorbed using $P(N,K)\ge1$.  This proves
the quantitative assertion of Theorem \ref{thm:quant}.

\section{The exponent requested in Question 6.1}

For $K>1$, put
\[
 a(K)=\frac{K-1}{4K} \in(0,1/4).
\]
On $0<a\le1/4$, one has $\gamma(a)\le4a$, and hence
$1+\gamma(a(K))<K$.  Since $P(N,K)$ increases with $K$, the theorem gives a
positive polynomial exponent at every $K>1$.

Nikolski's established estimate \cite{Nikolski}, quoted in the source as
equation (1.5), is
\[
 P(N,K)\ge\frac15N^{1-2/K}\qquad(K>\pi+1).
\]
Define
\[
 \alpha(K)=
 \begin{cases}
 \displaystyle\frac{a(K)\log2}{4(1+a(K)\log2)},&1<K\le\pi+1,\\[6pt]
 \displaystyle1-\frac2K,&K>\pi+1.
 \end{cases}
\]
The first branch is strictly increasing and tends to $0$ at $1^+$.  Its
value is less than $\log2/16$, whereas the right limit of the second branch
at $\pi+1$ is $1-2/(\pi+1)>1/2$.  Thus the piecewise function is strictly
increasing on $(1,\infty)$, and its limit at infinity is $1$.  The two lower
bounds prove $P(N,K)\gtrsim_KN^{\alpha(K)}$, completing the affirmative answer.

\section{Verification and limitations}

The script \texttt{code/verify\_tensor\_shift.py} checks the tensor product
identity, the scalar averaging inequality on finite grids, the power witness,
and sampled Kreiss ratios for three finite matrices.  All tests passed.  The
sampling is not a proof of the resolvent estimate; Proposition
\ref{prop:kreiss} is the proof.

The four lightweight run indexes and bounded arXiv/web searches were queried
using arXiv:2512.10025 and combinations of ``small Kreiss constants'',
``$P(N,K)$'', ``polynomial lower bound'', ``tensor power'', and ``weighted
shift''.  They returned the source and background material but no later answer
to Question 6.1.  The current arXiv version is v2 (11 March 2026), where the
question remains open.  This is not an exhaustive literature review; novelty
is unconfirmed.

The most important human-review points are Lemma \ref{lem:average}, the Abel
conversion in Proposition \ref{prop:kreiss}, and the exact-dimension padding.
Question 6.2 remains open.

\begin{thebibliography}{9}
\bibitem{CTY}
N. Chalmoukis, G. Tsikalas, and D. Yakubovich,
\emph{Operators with small Kreiss constants}, arXiv:2512.10025v2, 2026.

\bibitem{Nikolski}
N. Nikolski,
\emph{Sublinear dimension growth in the Kreiss matrix theorem},
Algebra i Analiz 25 (2013), no.~3, 3--51;
St. Petersburg Math. J. 25 (2014), 361--396.
\end{thebibliography}

\end{document}
